% Short running form; see the note on head width in format.tex.
\runninghead{David's Son and David's Lord}
\properlane{son-and-lord-of-david}{David's Son and David's Lord: The One God Confessed in His Word and Spirit}
\label{sec:son-and-lord-of-david}

Heard from Christ's own question, the Mass is a confession of faith. ``What
think you of Christ? Whose son is he?'' The Fathers answer that the Christ is
David's son according to the flesh and David's Lord according to his eternal
birth from the Father, equal to him in honour and power. The Epistle's ``one
Lord, one faith, one baptism, one God and Father of all'' is read by some of
them as the confession of the Trinity, and by others as the one God's lordship,
providence and indwelling, which the Son shares without inferiority. The
Gradual's ``By the word of the Lord the heavens were established, and all the
power of them by the spirit of his mouth'' is the verse in which the Fathers
hear the Father creating by his Word and his Spirit. The God whom the Collect
asks to follow ``alone'', and whom the Gospel commands to be loved with the
whole heart, is then the God whose Son sits at his right hand; and the
Preface appointed for this Sunday confesses him as ``one God, and one Lord
\dots\ in a Trinity of one substance''.

\subsection{Christ's question and the Fathers' answer}

Jesus asks the Pharisees whose son the Christ is. They answer, ``David's''.
He replies with Ps~109:1, ``The Lord said to my Lord: Sit on my right hand,
until I make thy enemies thy footstool'', and asks, ``If David then call him
Lord, how is he his son?'' The point turns on David's speaking ``in spirit'':
the psalmist calls his own descendant his Lord, and he does so by the Spirit.

St Augustine opens his exposition of Ps~109 with this Gospel. The psalm
prophesies Christ \latin{certe aperteque}, surely and openly. Christians must
answer Christ's question as the Pharisees did, ``David's'', \latin{sed non
remaneamus ubi Judaei}: but they must not stay where the questioners stayed.
The answer the Pharisees could not give is the prologue of John, which
Augustine addresses to Christ himself. Thou wast in the beginning the Word,
and all things were made by thee: \latin{ecce Dominus David}, behold David's
Lord. For our weakness the Word was made flesh to dwell in us: \latin{ecce
filius David}, behold David's son. Or, in the words of Philippians, being in
the form of God he is David's Lord, and emptying himself, taking the form of a
servant, he is David's son. In asking ``how is he his son?'', Augustine tells
Christ, \latin{non te filium ejus negasti, sed modum quo id fieret
inquisisti}: he did not deny that he was David's son, but asked how it came to
be. As for the questioners, \latin{maluerunt inflata taciturnitate disrumpi,
quam humili confessione edoceri}: they preferred to burst with a swollen
silence rather than be taught by a humble confession (\work{Enarrationes in
Psalmos} 109, 3--6).

St Jerome presses the same logic from the side of the psalm. If the Christ is
a mere man and only David's son, how can David call him his Lord, and that not
by error or by his own will but in the Holy Spirit? \latin{Dominus igitur
David vocatur, non secundum id quod de eo natus est, sed juxta id quod natus ex
Patre semper fuit, praeveniens ipsum carnis suae patrem}: he is called David's
Lord not by his birth from David but by his birth from the Father, which always
was, and which goes before the father of his flesh.
Jerome meets a Jewish reading that applied the psalm to Abraham with the
psalm's own later words, ``from the womb before the day star I begot thee''
and ``a priest for ever'' (\work{Commentarii in Matthaeum} IV, on 22:41--46).

St John Chrysostom attends to the way Christ teaches. Jesus had just praised
the confession that there is one God, and he explains why that praise does not
exclude himself: the saying ``There is one God'' is spoken ``not to the
rejection of the Son, but to make the distinction from idols''. The
counter-question follows, ``secretly leading them on to confess Him also to be
God''. Since they held him a mere man, he brings in David witnessing ``to His
being Lord, and the genuineness of His Sonship, and His equality in honor with
His Father''. He adds the second half of the verse, ``Till I make Thine
enemies Thy footstool'', ``in order to move them to fear''. He teaches by
question and inference, ``in condescension'', so that the saying might not give
offence (\work{Homilies on Matthew} 71).

St Thomas Aquinas gathers these threads into a doctrinal statement. Christ has
two generations, and the Pharisees answered only of the flesh. The psalm shows
three things of the Son: \latin{praeeminentiam ad sanctos, aequalitatem ad
patrem, et dominium super rebelles}, his pre-eminence over the saints, his
equality with the Father and his dominion over the rebellious. The conclusion
is plain: \latin{filius est secundum carnem \dots\ et dominus secundum
divinitatem} (\work{Super Matthaeum} c.~22, lect.~4). The liturgical
commentator William Durandus, who knew churches that read this Gospel on the
following Sunday, reads its second half as Christ reproving the Pharisees
\latin{de infidelitate, quia non credebant ipsum esse Deum}: for their
unbelief, because they did not believe him to be God (\work{Rationale} VI, on
the Eighteenth Sunday).

The three Fathers and the Doctor agree on the answer to Christ's question. They
differ in what they bring out: Augustine the form of God and the form of the
servant, and the humble confession that the proud silence refused; Jerome the
eternal birth that precedes the father of the flesh; Chrysostom the patient
pedagogy that leads toward confession and the equality of honour it reveals;
Aquinas the threefold dignity of the Son.

\subsection{One Lord, one faith, one baptism, one God and Father}

Paul grounds the Church's unity in a sevenfold ``one'' that ends with ``One
God and Father of all, who is above all, and through all, and in us all''.
The Fathers who comment on it read the last phrases in two ways.

St Jerome, writing against Sabellius and the Arians, gives the three phrases
to the three Persons. \latin{Unus est Dominus, et unus est Deus, quia Patris
et Filii dominatio, una divinitas est}: one Lord and one God, because the
lordship of Father and Son is one godhead. The faith is one because we
believe alike in Father, Son and Holy Spirit; the baptism is one, though we are
immersed three times, \latin{ut Trinitatis unum appareat sacramentum}, so that
the one mystery of the Trinity may be shown. Then the phrases themselves:
\latin{Super omnes enim est Deus Pater, quia auctor est omnium. Per omnes
Filius, quia cuncta transcurrit, vaditque per omnia. In omnibus Spiritus
sanctus, quia nihil absque eo est}: the Father is above all as author of all,
the Son through all as passing through and pervading all things, the Holy
Spirit in all, because nothing is without him (\work{Commentarii in epistulam
ad Ephesios} II, on 4:4--6).

St Thomas Aquinas reaches the same end by two steps. He reads ``One Lord'' as
Christ, ``one faith'' as both the thing believed and the habit of believing,
and ``one baptism'' as one because Christ baptizes within, because it is
given \latin{in nomine unius, scilicet trinitatis}, and because it is not
repeated. ``One God'' recalls \latin{Audi Israel, Dominus Deus tuus unus est}, and
``Father of all'' God's kindness toward us. Then he commends the one God's
dignity \latin{ex tribus}: the height of his godhead, ``who is above all''; the
breadth of his power, ``through all''; the largesse of his grace, ``in us
all''. Only after reading the phrases of the one God does he appropriate them:
\latin{Sed primum appropriatur patri qui est fontale principium divinitatis},
the first to the Father as the fountain-source of the godhead, the second to
the Son, who is wisdom reaching from end to end mightily, the third to the
Holy Spirit, who fills the whole world (\work{Super Ephesios} c.~4,
lect.~2).

St John Chrysostom reads the verse otherwise. In his eleventh homily on the
letter, in the English of the Nicene and Post-Nicene Fathers: ```Who is over
all,' that is, the Lord and above all; and `through all,' that is, providing
for, ordering all; and `in you all,' that is, who dwelleth in you all. Now this
they own to be an attribute of the Son; so that were it an argument of
inferiority, it never would have been said of the Father.'' He does not give
the three phrases to three Persons. He reads them as three relations of the
one God and Father to all things: over all as Lord, providing for all,
dwelling in all. Theodoret of Cyrus names the same three relations lordship,
providence and indwelling. What Chrysostom adds is an argument about the Son:
an attribute that is also ascribed to the Son, and is here said of the Father,
cannot mark inferiority in the Son.

These are two lines of exegesis, and neither denies the other. The
distribution of ``above all, through all, in all'' among the Persons rests on
Jerome and on Aquinas's appropriation; the Greek commentators read the phrases
of the one God and Father, and Chrysostom draws from them the Son's equality.
Aquinas holds both, reading first of the one God and then appropriating. For
the confession this reading follows, Jerome and Aquinas supply the Trinitarian
sense of the verse, and Chrysostom supplies its bearing on the Son. The Epistle
thus confesses what the Gospel's question uncovers: the one Lord whom Paul
names is the Christ whom David calls Lord.

\subsection{By the word of the Lord and the spirit of his mouth}

The Gradual's versicle, Ps~32:6, gathers a remarkable consensus.
St Irenaeus of Lyons, stating the rule of faith against those who divided God,
uses it as a proof text: ``The rule of truth which we hold, is, that there is
one God Almighty, who made all things by His Word \dots\ Thus saith the
Scripture \dots\ `By the Word of the Lord were the heavens established, and
all the might of them, by the spirit of His mouth.'\,'' God ``by His Word and
Spirit, makes, and disposes, and governs all things'', and is ``the God of
Abraham, and the God of Isaac, and the God of Jacob, above whom there is no
other God'' (\work{Adversus haereses} I.22.1). Irenaeus does not expound the
psalm; he uses its verse to name the one Creator who works by his Word and his
Spirit.

St Basil the Great, in his treatise on the Holy Spirit, reads the verse's
terms directly: ``the Word is He who was with God in the beginning and was
God, and the Spirit of the mouth of God is the Spirit of truth which proceeds
from the Father. You are therefore to perceive three, the Lord who gives the
order, the Word who creates, and the Spirit who confirms'' (\work{De Spiritu
Sancto} 16.38). St Augustine, preaching on the psalm, stops his hearers at the
same verse: \latin{Sane, fratres, videte eadem opera Filii et Spiritus sancti
\dots\ Verbum certe Dei Filius est, et Spiritus oris ejus Spiritus sanctus
est}: the works are the same works of the Son and of the Holy Spirit, for the
Word of God is the Son and the Spirit of his mouth is the Holy Spirit. And
lest the Father be left out: \latin{Hoc ergo facit Filius et Spiritus sanctus.
Numquid sine Patre?} The Son and the Spirit do this; but not without the
Father (\work{Enarrationes in Psalmos} 32, II, s.~2, 5). Cassiodorus finds the Trinity in the grammar of the verse itself.
In saying \latin{Verbo} the psalm declares the Son, in adding \latin{Domini}
the Father, in \latin{spiritu oris eius} the Holy Spirit, \latin{et ut in
tribus personis manifestam intelligeres unitatem, eius dixit, non eorum}: and
to make the unity in three persons plain, the psalm says \latin{eius}, his,
and not \latin{eorum}, theirs (\work{Expositio psalmorum}, on Ps~32:6).

Aquinas reads the verse by appropriation, as he reads Eph 4:6.
\latin{Dominus est nomen potestatis, et potentia appropriatur patri}; the Word
is the Son; the Spirit of his mouth is the Holy Spirit, ``spirit of the mouth''
because the mouth is appropriated to the Word, so that it means the spirit of
the Word. The works of the Trinity are undivided, but the psalm speaks
\latin{secundum appropriationem} (\work{Super Psalmos} 32, n.~5). St Robert
Bellarmine keeps the literal sense, God's word of command in the first chapter
of Genesis, and the mystery together. Although God's orders are what the verse
means, ``there is no doubt but the Holy Ghost meant to glance at the mystery of
the Holy Trinity to be revealed in the New Testament''. He guards the unity at
once: ``the acts of the Trinity cannot be separated, by reason of the unity of
essence''. Augustine, Cassiodorus and Aquinas also read the ``heavens'' made
firm as the apostles, weak men strengthened by the Word and the Spirit.

The Missal prints the versicle at this Sunday with a small letter,
\latin{spíritu}, as it prints the \latin{spíritus} of the Epistle; setting the
same Gradual in the votive Mass of the Holy Spirit, it prints \latin{Spíritu}.
Basil, Augustine and Cassiodorus read the Word and the Spirit out of the
psalm's own words.

\subsection{The other texts as confession}

The Collect's \latin{te solum Deum} asks to follow the one God alone. Heard
beside the Gospel, it is Chrysostom's point in prayer: ``one God'' shuts out
the idols, not the Son. Its earliest witness reads \latin{te solum Dominum};
both words stand together in the Gospel's \latin{Diliges Dominum Deum tuum}.

The Alleluia divides its witnesses here as elsewhere. For Augustine the poor
man who prays is the eternal Word made poor for us. \latin{Unde illi labor,
unde gemitus Verbo, per quod facta sunt omnia?}, he asks: whence labour and
groaning to the Word through whom all things were made? He answers that the
Word deigned to take our death. In the same place he sets three unions side by
side and distinguishes them by the gender of a single pronoun:
\latin{Christus et Ecclesia, utrumque unus: sed Verbum et caro non utrumque
unum; Pater et Verbum utrumque unum; Christus et Ecclesia utrumque unus, unus
quidam vir perfectus}. Christ and the Church are \latin{unus}, masculine: the
two are one man, the perfect man of Eph 4:13. The Father and the Word are
\latin{unum}, neuter: one thing, one in substance. The Word and the flesh are
not \latin{unum} in that sense, since Godhead and manhood are not one
substance; they are one person, and only because they are one person can the
head be called \latin{unus} with his body in the same breath
(\work{Enarrationes in Psalmos} 101, s.~1, 2). That distinction is what this
reading needs. The one Christ is David's son in the flesh he took and David's
Lord in the substance he has of the Father, and the second confession takes
nothing from the first. Cassiodorus, by contrast,
declines to put the psalm in Christ's mouth and hears the prophet speaking for
the poor of Christ. Neither Father yields to the other; only Augustine's
reading belongs to the confession of the Word.

The Offertory reaches Christ through the answer Daniel received. While he was
still praying, Gabriel came to him (Dan 9:20--21), and Jerome remarks
\latin{grandis orationis effectus}: great was the prayer's effect. The
angel brought the prophecy of the seventy weeks (9:24--27). Jerome sets out the
computations of Africanus, Eusebius, Hippolytus, Apollinaris, Clement,
Tertullian and the Jews, and declines to choose among them:
\latin{periculosum est de magistrorum Ecclesiae judicare sententiis}, since it
is perilous to sit in judgment on the opinions of the Church's masters
(\work{Commentarii in Danielem}, on 9:21--24). Rupert of Deutz, expounding the
chant with its old verse \latin{Adhuc me loquente}, says that Daniel was so
exalted by the vision that he not only foresaw Christ coming for our salvation
but learnt beforehand the times of his coming (\work{De divinis officiis}
XII.17). Durandus likewise says that the archangel made the incarnation of
Christ known to him (\work{Rationale} VI.131). The 1962 antiphon ends at the
petitions of 9:17--19 and prints no verses; the answer that the commentators
join to the prayer belongs to the chapter that follows it.

Read beside the Gospel, the Communion's God ``who taketh away the spirit of
princes'' and is ``terrible with the kings of the earth'' is the Lord of
Ps~109, to whom every enemy is made a footstool. In the Gospel's psalm Aquinas
finds the Son's \latin{dominium super rebelles}, and Chrysostom says the
footstool is added ``to move them to fear''. The
Introit's \latin{Iustus es, Domine} is addressed to that one Lord; Cassiodorus
hears in it the confession by which Daniel was saved among the lions, a link to
the Offertory's prophet that is his own.

The Postcommunion adds nothing distinctive to this reading. Its address is to
the \latin{omnípotens Deus}, with no distinction of Persons; what it asks ---
that by God's \latin{sanctificationes} our vices be cured and eternal remedies
come to us --- is the healing the first and the third readings develop, not a
confession of the Word and the Spirit; and its only Trinitarian words are the
conclusion it shares with the Collect and the Secret, which the typical
edition prints here merely as \latin{Per Dóminum nostrum} and which the general
rubrics complete (no.~115~a).

The Secret addresses the Father's \latin{maiestas}. Its words name neither the
Son nor the Spirit, and it concludes, like the other two orations, through the
Son who lives and reigns with the Father in the unity of the Holy Spirit. Its
first word meets the Preface the Missal directs immediately after it, whose
last clause adores \latin{in maiestáte \dots\ æquálitas}. That adjacency belongs
to every Sunday of the second class that uses this Preface. The Preface itself
is the confession this Mass makes, and it gathers what the Gospel and the
Epistle have shown. In the 1861 English: ``Who together with thy only begotten
Son and the Holy Ghost, art one God, and one Lord not in a singularity of one
Person, but in a Trinity of one substance \dots\ So that in the confession of
the true and eternal Deity, we adore a distinction in the Persons, an unity in
the essence, and an equality in the Majesty.''

\subsection{The difficulties this confession must carry}

Three texts resist being made simple witnesses of the Trinity. The
distribution of Eph 4:6 among the Persons rests on Jerome and on Aquinas's
appropriation; Chrysostom and Theodoret read the verse of the one God and
Father, and Chrysostom's support for the Son comes by argument, not by
distribution. The Christological Alleluia is Augustine's and is contested by
Cassiodorus. The Offertory's Christological weight lies in verses that the
1962 antiphon does not sing; the older verses and the medieval commentators
supply it. The Preface, finally, is the ordinary Sunday Preface of the season
and speaks for every such Sunday. What remains firm is the Gospel's second half,
on which the Fathers are unanimous, and the Gradual's versicle, on which
Irenaeus, Basil, Augustine, Cassiodorus, Aquinas and Bellarmine agree in
hearing the one God who creates by his Word and his Spirit.

\subsection{The four senses of this reading}

\begin{fourSenses}
\item[Literal.] Jesus asks the Pharisees whose son the Christ is and, when
they answer ``David's'', asks how David, speaking by the Spirit in Ps~109, can
call him Lord; no one can answer him. Paul grounds the Church's unity in one
Lord, one faith, one baptism and one God and Father of all. The psalmist praises
the Lord by whose word the heavens were made firm and all their host by the
breath of his mouth.

\item[Allegorical.] Christ is David's son in the flesh and David's Lord from
eternity, in the form of the servant and in the form of God (Augustine,
Jerome, Chrysostom, Aquinas). The one God who creates by his Word and his
Spirit is Father, Son and Holy Spirit (Irenaeus, Basil, Augustine,
Cassiodorus, Aquinas, Bellarmine), above all, through all and in all (Jerome,
Aquinas). The heavens made firm by the Word and the Spirit are also the
apostles made strong (Augustine, Cassiodorus, Aquinas).

\item[Moral.] Answer Christ's question as the Church does, and do not stop
where his questioners stopped; exchange a proud silence for a humble
confession (Augustine). Guard the one faith and the one baptism against the
division that breaks the bond of peace (Jerome). Follow the one God alone,
knowing that ``one God'' shuts out the idols and not the Son (Chrysostom).

\item[Anagogical.] The Son sits at the Father's right hand until every enemy
is made his footstool (Ps~109:1), with dominion over the rebellious (Aquinas).
The blessed nation of the Gradual's psalm is, for Bellarmine, most happy
``when we shall see him as he is'', and for Aquinas its beatitude, begun now in
hope, is to cleave to God by knowledge and love, perfectly hereafter. The
vision so promised is of the God whom the Preface confesses, adored in a
distinction of Persons, a unity of essence and an equality of majesty.
\end{fourSenses}
